\documentclass[]{../../../../tex/classes/styledArticle}
\usepackage{../../../../tex/packages/hebrewSupport}


\author{שחר פרץ}
\title{\textit{היסטוריה 48}}
\begin{document}
	\maketitle
	\subsubsection*{מבנה המבחן}
	תזכורת נוספת לגבי מבנה המבחן: בבחינה ארבעה פרקים.
	\begin{itemize}
		\item פרק א' חובה, טוטאליטריות ושואה, 3 שאלות (לרוב עם מקור, ייתכן שאחת ללא). יש לענות על שאלה אחת מתוכן. 
		\item פרק ב', לאומיות וציונות. 
		\item פרק ג', בונים מדינה. 
		\item פרק ד', סוגיות בתולדות מדינת ישראל. 
	\end{itemize}
	בכל אחד מהפרקים 3 שאלות ובוחרים אחת. שרית שוב ממליצה שלא להמר וללמוד את כל החומר. בפרק ד' לא למדנו עלייה וקליטה, בפרק ג' לא למדנו את כל מלחמת העצמאות. בפרקים הראשונים למדנו הכל. הערות: 
	\begin{itemize}
		\item לוודא שיש ספרים לכולם
		\item לא להעתיק פסקאות שלמות
	\end{itemize}
	
	במתכונת, בפרק ד' לא יהיו 3 שאלות, אלא רק 2 עליהן נדע לענות. בפרק ג' יהיו 3 שאלות אבל לא בהכרח נדע לענות על כולן. 
	
	\subsubsection*{עניה על שאלה}
	בוחרים שאלה, מסמנים מילות הוראה, מומלץ לעשות שלד של תשובה. אם זה מקור, קראו וחפשו את התשובה. התחילו במשפט פתיחה (לא עמוד וחצי פתיחה). ציטוטים הם לא תשובה. ציין הצג הסבר. משהו. שאלות בלי הוראה או עם הוראה לא ברור, זה הסבר. מומלץ לכתוב ברמת הסבר תמיד, לא יורידו על זה נקודות. אין מספר שקול לשלושה פריטים. אם צריך 2 או 3 ממשהו, בודקים את הראשונים. אם כתוב משהו כמו ''הציגו שלושה לפי המקור ולפי מה שלמדתם``, צריך לפחות אחד מהמקור ולפחות אחד שלא מופיע במקור. 
	
	בשאלות עמדה, עמדה כוללת טענה ונימוק, כאשר הנימוק מובסס על עובדות היסטוריות. עובדות היסטוריות=מה שלמדתם. לקחת אירועים, לא לקחת את תנועת המרי כאשר שואלים מה השפיע יותר על הבריחה של הבריטים. 
	
	בשאלות בהן מוצג קטע מקור וכותבים למעלה תאריך ושנה, יש לכך חשיבות. למטה, במקור של קטע המקור, היא כתוב ''מתוך ספר של בלי בלה בלי בלה`` – זה לא רלוונטי לשאלה (אולי צריך את זה להצגה של קטע המקור). 
	
	בשאלות השוואה – תבחין, תבחין תבחין תבחין בשורות שונות, כל אחד מהם עושים דומה (מאוחד) ושונה. בעמודות מקור א' מקור ב'. אל תעשו השוואה בראש. 
	
	\subsubsection*{מאפייני מקור}
	להתייחס למאפיינים הבאים: מי (כתב/אייר/צייר), מתי נכתב המקור, היכן נכתב המקור, עבור מי, וכו'. במקורות חזותיים לפרט. 
	
	\subsection*{תבניות}
	\subsubsection*{תבניות לקטעי מקור}
	לפניכם שני קטעי מקור העוסקים ב־??? / שנוצרו בתקופת ???. מקור 1 הוא ??? מקור 2 הוא ???. הציגו הבדל/הבדלים בין שני קטעי המקור בנוגע ל???/בתיאור של ???. הסבירו 2 סיבות להבדל. בתשובתכם התייחסו למאפייני המקור או להקשר ההיסטורי בו הוא נוצר. 
	
	לפניכם קטע מקור שעוסק ב???. [משהו מבין הדברים הבאים, יכול להיות שניים מהם] הציגו את הטיעון (טענות ונימוקים) של יוצר המקור. הציגו טענה נגדית ונמקו אותה באמצעות שתי עובדות היסטוריות. הציגו נימוקים נוספים לחיזוק הטענה שבמקור ובססו אותם באמצעות שתי עובדות היסטוריות / עובדה ומקור / שני מקורות. 
	
	לפניכם קטע מקור (מאפייני המקור ???) העוסק ב???. הציגו את העמדה/התיאור בנוגע ל??? שבאה לידי ביטוי בקטע המקור. הסבירו את השפעת מאפייני המקור/ההקשר ההיסטורי שבו נוצר המקור/התיאור שהצגתם. הציגוג את העקרונות/תפיסות/רעיונות/בנוגע ל??? הבאים לידי ביטוי בקטע המקור. תארו פעולה לקידום/מימוש כל אחד מהעקרונות/תפיסות/רעיונות שהצגתם, והסבירו כיצד כל אחת מן הפעולות שתיאםתם קידמה/ממשה עקרונות/תפיסות/רעיונות. 
	
	קיימות אפשרויות נוספות. 
	
	\subsubsection*{תבניות למקור}
	הציגו את עמדתכם/העמדה/עמדות של ??? בויכוח/דילמה/דיון בנושא ???. בחרו מקור/ות מספר הלימוד ונמקו באמצעות המקור את העמדה/עמדות שהצגתם. 
	
	חוקרים/היסטוריונים טוענים ש''???'' (היגד/ציטוט/טענה). הסבירו את הטענה של ???. בחרו מקורות/ות מספר הלימוד ונקמרו באמצעותו/ם את הטענה שהסברתם. בחרו מספר הלימוד ונמקו באמצעותו את הטענה הנגדית. 
	
	הציגו שניים/שלושה מן העקרונות/מאפיינים/תוצאות/רעיונות של ???. בחרו מקור/ות והראו כיצד שניים/שלושה מהמפיינים באים לידי ביטוי במקור. 
	
	\subsubsection*{שאלות ללא מקור}
	לפניכם שני דפוסים/שלבים של/תגובות ל???, הציגו נקודת דמיון אחת/הבדל/ים בינהם על פי התבחין/ים, הסבירו את הגורם/גורמים להבדלים שהצגתם. הסבירו מה ניתן ללמוד על (מדיניות/מטרות) משלושת הפעולות/צעדים שנעשו בתחום/ים ב??? (מקום, שנים). הציגו שלושה השקפות/גישות של ??? בנוגע ל??? והסבירו את הסיבה/ות שהובילה/ו להבדלים אלה. הציגו שתיים מן המטרות/הקשיים של ??? בתקופת ה??? והסבירו באמצעות שתי פעולות/צעדים כיצד הושגו המטרות/כיצד התמודדו עם הקשיים. הציגו שניים מן השינויים שהתרחשו ב(נושא/מקום) בשנים ??? והסבירו כיצד טרם כל אחר מן השינויים למימוש/קידום/צמיחה. 
	
	\subsubsection*{שאלות מבט על}
	שאלות עם כותרת כללית כמו ''WWII בתקופה דכגידגחכל``. 
	
	הציגו את האירוע/המפיין/הפעולה/התוצאה/האישיות אשר לטענכם הוא החשוב/משמעותי ביותר בתהליך/בתקופת, בססו את טענתכם באמצעות שני נימוקים. הציגו אירוע המסמל את נקודת ההתחלה של תהליך/תקופה, ואירוע המווסה נקודת סיום של תהליך/תקופה זו, נמקו את הבחירה בכל אחד מהאירועים באמצעות עובדות היסטוריות. הציגו אחד/שניים מן השינויים שחלו בנקודת המבט/בתפיסה/בגישה של ??? בנוגע ל??? בשנים/העשורים ??? והסבירו את הגורמים לשינוי/ים שהצגתם. לפניכם רשימה של אירועים/מושגים/דמויות היסטוריות, בחרו שלושה מהם, הציגו'הגדירו/תארו כל אחד מהם והסבירו את הקשר בינהם. (היגד/קביעה המבטא על התקופה ???), הצידו שלוש עובדות היסטוריות... 
	
	לדוגמה: תבנית היכד/קביה המבטא מבט על תקופה, הציגו 3 עובותד מנקודות זמן שונות ונמקו באמצעותן את ההיגד/קביעה. נתבונן בשאלה ''בכל אחד משני השלבים שבמלחמת האמצעות היו תקופות שבהן עבד הצד היהודי מדרך פעולה של מגננה לדרך פעולה של יוזמה והתקפה. הסבירו טענה זאת, ונמקו אותה באמצעות שתי עובדות היסטוריות לכל שלב במלחמה``. 
	
	זו שאלת מבט על. אין בעיה לעשות אותה למרות שבקושי למדנו על מלחמת העמצאות. יש שלב א' שלב ב', בכל אחד מהם יש מעבר ממגננה ליוזמה והתקפה. נמצא אירועים ונסביר איך הם קשורים, לדוגמה להראות שמבצע נחשון הוא מתקפה יזומה ושינוי בתפיסה. 
	
	לדוגמה: תבנית הסיטוריונים טוענים???ף הסיברו את הטענה של???. נתבונן בשאלה ''היסטוריונים טוענים שהפתרון הסופי לא תוכנן מראש. נמקרו טענה זו באמצעות 3 עובדות היסטוריות``. 
	
	
	\subsection*{דגשים}
	ליהודים אין מקום בסדר החדש, לכן הגטאות לא חלק ממימוש הסדר החדש. אם מסבירים תמותה טבעית וזה עוד אפשר לקבל חלקית. בסדר החדש אין יהודים. לכן מחנות השמדה מתקבלים כמימוש הסדר החדש, אך גטאות ומחנות ריכוז לא. 
	
	תיקון: במתכונת תהיה שאלה של עלייה וקליטה או ששת הימים, ויהיה כתוב מפורשות שלא למדנו. 
	
	
\end{document}